% V-model mapping of the two-layer verification strategy.
% Kept as diagram source rather than a bitmap so that it diffs, reviews and
% rebuilds deterministically with the rest of the document.
\begin{tikzpicture}[
    x=1mm, y=-1mm,
    phase/.style={
      draw, fill=black!5,
      text width=38mm, align=center,
      inner sep=1.6mm, font=\small\bfseries,
    },
    stage/.style={
      draw, text width=38mm, align=center,
      inner sep=1.6mm, minimum height=14mm, font=\footnotesize,
    },
    flow/.style={-{Latex[length=2mm,width=1.6mm]}, thick},
    corr/.style={-{Latex[length=2mm,width=1.6mm]}, thick, densely dashed},
  ]

  \node[phase] (specphase) at (24, 6) {Design \& Spezifikation};
  \node[phase] (vvphase)   at (118, 6) {Verifikation \& Validierung};

  \node[stage] (req)  at (24, 34) {\textbf{Systemanforderungen}\\(Legacy-Verhalten nach Originalhandbuch)};
  \node[stage] (arch) at (36, 62) {\textbf{Architektur-Design}\\(FPGA/STM32-Trennung, Modul-Aufbau)};
  \node[stage] (mod)  at (48, 90) {\textbf{Modul- \& Komponenten-Design}\\(VHDL-Logik, Schaltpläne)};

  \node[stage] (unit)  at (106, 90) {\textbf{Unit- \& Modultests}\\\textbf{(Layer 1)}\\(VHDL-Simulation)};
  \node[stage] (integ) at (118, 62) {\textbf{Integrationstests}\\\textbf{(Layer 1)}\\(Hardware-Messungen, Schnittstellen)};
  \node[stage] (sys)   at (130, 34) {\textbf{Systemvalidierung}\\\textbf{(Layer 2)}\\(Fixture-Tests, Szenarien)};

  \draw[flow] (req)   -- (arch);
  \draw[flow] (arch)  -- (mod);
  \draw[flow] (mod)   -- (unit);
  \draw[flow] (unit)  -- (integ);
  \draw[flow] (integ) -- (sys);
  \draw[corr] (arch)  -- (integ);

\end{tikzpicture}
